\section{Lesson 2, 29/09}
\subsection{Internal Energy as a TD potential}
[10 min. cfr with notes pt 2 on moodle/Remi]
$N = const.$
\subsection{Legendre Transform}
What if the control variables are the intensive ones, e.g. system in a thermostat? It
is useful to use $T$ instead of $S$: the mathematical tool used is the \textit{Legendre
transform} (the one used to transform $L(q,\dot q)\rightarrow H(q,p)$, using 
$p = \pdv[]{L}{\dot q}$)
$$E(S,V,N,X)\rightarrow F(T,V,N,X):=E\left(S,V,N,X\right)-TS$$
where $S$ is found by inverting the relation $T=\pdv[]{E}{S}$. F is the Helmholtz
free energy

Using the fundamental equation of TD
$$E=TS-pV+\mu N+yX$$
the function $F$ can also be written as
$$F = -pV +\mu N +yX:$$
differentiating this yields
$$\dd F \leq -S\dd T-p\dd V + \mu \dd N + y\dd X$$
so keeping fixed $T,V,N,X$, the equilibrium is at minimal $F$.

What are the conjugated variables? As always
\begin{align*}
S = -\pdv[]{F}{T}\\
p = - \pdv[]{F}{V}\\
\mu = \pdv[]{F}{N}
\end{align*}
This process can be done with every intensive variable e.g. enthalpy:
$$E(S,V,N,X)\underset{p = -\pdv[]{E}{V}}\longrightarrow H(S,p,N,X) := E(S,V,N,X) +pV$$
and
$$H=TS+\mu N+yX \underset{\dd\text{ and fund. eq.}}\longrightarrow 
\dd H \leq T\dd S + v \dd p +\mu \dd N +y\dd X$$
so fixing $S,p,N,X$, the equilibrium is at minimal $H$.

Doing this with respect to both couples $(T,S)$ and $(p,V)$ yields the Gibbs free
energy:
$$G(T,p,N,X) = E(S,V,N,X) - TS + pV$$
using the fundamental equation yields
$$G = \mu N + yX$$
\textbf{Note} In simple systems (without $(y,X)$) $G=\mu N$. The differential is
$$\dd G\leq -S\dd T + V\dd p +\mu \dd N +y\dd X$$
that shows that, fixing $T,p,N,X$, equilibrium is at minimal $G$.

The last Legendre Transform of interest is the one that yields the \textit{Grand
Canonical potential}:
$$E(S,V,N,X)\underset{(T,S),(\mu,N)}\longrightarrow \bar\omega(T,V,\mu,X) = E(S,V,N,X)-
TS-\mu N$$
The fundamental equation yields
$$\bar\omega = -pV+yX$$
and differentiating yields
$$\dd \bar\omega\leq -S\dd T -p\dd V -N\dd\mu+y\dd X$$
so the equilibrium is at the minimum of $\bar\omega$.

The pressure is simply
$$ p = -\pdv[]{\bar\omega}{V} = -\frac{\bar\omega}{V}$$
\subsection{Response Functions}
These functions describe the reaction of variables when another one changes (typically
conjugated variables), are very easy to measure and they are deeply connected to the
size of the fluctuations of the system (Fluctuation-dissipation theorem).

\textbf{Def.} Heat capacity
$$C = \frac{\delta Q}{\dd T}$$
is the amount of heat to be furnished to change the $T$; it is not an exact
differential, so it depends on the transformation.\\
Heat capacity at constant $(V,N)$: if $\dd V = \dd N = \dd X = 0 \rightarrow\delta W=0$;
by the first principle
$$\delta Q = \dd E$$
so
$$C_{V,N,X} = C_V = \frac{\delta Q}{\dd T} = \eval{\pdv[]{E}{T}}_{V,N,X}$$
or, using the definition of entropy
$$\delta Q = T\dd S$$
and fixing $(V,N,X)$
$$\dd S = \eval{\pdv[]{S}{T}}_{V,N,X}\dd T$$
so
$$C_V = \frac{T\dd S}{\dd T} = T \left(\pdv[]{S}{T}\right)_{V,N,X}.$$

A more natural approach springs from the Helmholtz free energy:
$$S = -\eval{\pdv[]{F}{S}}_{V,N,X}\rightarrow C_V = - T \pdv[2]{F}{T}$$

Heat capacity at constant $(p,N,X)$: the first principle states
$$\delta W = p \dd V \rightarrow \delta Q = \dd E +p\dd V$$
so
$$C_p = \frac{\delta Q}{\dd T} = \eval{\pdv[]{E}{T}}_{p,V,X} +
p\eval{\pdv[]{V}{T}}_{p,V,X}$$
but also
$$E = G +TS-pV = G - T\pdv[]{G}{T}-pV$$
so
$$C_p = \pdv[]{}{T}\left(G-T\pdv[]{G}{T} -pV\right) + p\pdv[]{V}{T} = -
\eval{T\pdv[2]{G}{T}}_{p,N,X}$$
\subsection{Mechanical response functions}
Susceptibility $\chi$ is the volume response to a change in pressure: this also varies
depending on the transformation.

Isothermal: fixing $(T,N,X)$
$$\chi_T := -\eval{\pdv[]{V}{p}}_{T,N,X}$$
since the variables are $(T,p,N,X)$, with only p independent, the most natural
TD potential to use is $G$: some calculations show that
$$\chi_T = - \eval{\pdv[2]{G}{p}}_{T,N,X}$$

Adiabatic: fixing $(S,N,X)$
$$\chi_S = -\eval{\pdv[]{V}{p}}_{S,N,X}$$
the TD potential is the enthalpy $H$:
$$V = \eval{\pdv[]{H}{p}}_{S,N,X} \rightarrow \chi_S = \eval{\pdv[2]{H}{p}}_{S,N,X}$$

All these response functions are between conjugated variables: what happens if they are
not? Consider the Thermal expansion coefficient
$$\alpha_{p,N,X} = \alpha_p = \eval{\pdv[]{V}{T}}_{p,N,X} \underset{\text{exercise}}=
\eval{\pdv[2]{G}{p}{T}}_{p,N,X}$$

\subsection{Maxwell Relations}
These are relations between TD quantities: consider the set of independent variables
$$\{Z_A\}$$
their conjugated variables
$$\{\xi_A\}$$
and a TD potential
$$\dd\Phi = \sum_{}^{}\xi_A\dd Z_A$$
it follows that
$$\pdv[]{\xi_A}{Z_B} = \pdv[2]{\Phi}{Z_A}{Z_B} = \pdv[2]{\Phi}{Z_B}{Z_A} =
\pdv[]{\xi_B}{Z_A}$$
from other mathematical properties (like the "chain rule", see Billò notes) arise other
relations. One finds
$$\frac{1}{T}C_p > \frac{1}{T}C_V + \frac{\alpha^2_p}{\chi_T}$$

\subsection{Stability of equilibrium state}
Even at equilibrium, at maximum $S$, there can be small fluctuations in TD variables.
For all such fluctuations, the entropy \textbf{must decrease} to ensure stability.

What are the conditions for local equilibrium? Consider an isolated system at
equilibrium subdivided in 2 parts, which can exchange particles and heat.
\begin{gather*}
E=E_1+E_2\\
V=V_1+V_2\\
N = N_1+N_2
\end{gather*}
Suppose some fluctuation occurs in $E_i,V_i,N_i$: the other part of the system must
have an opposite fluctuation
$$\Delta E_2 = -\Delta E_1 ...$$
Assuming small fluctuations, the variation of $S$ at first order is
$$\Delta S = \sum_{i=i}^{2} \left(\eval{\pdv[]{S_i}{E_i}}_\text{eq}\Delta E_i +
\eval{\pdv[]{S_i}{V_i}}_\text{eq}\Delta V_i + \eval{\pdv[]{S_i}{N_i}}_\text{eq}
\Delta N_i \right)=\left(\frac{1}{T_1^\text{eq}}- \frac{1}{T^\text{eq}_2}
\right)\Delta E_1 + \left(\frac{p_1^\text{eq}}{T_1^\text{eq}}-
    \frac{p_2^\text{eq}}{T^\text{eq}_2}\right) + ...$$
Since $\Delta S \leq 0$, it must be that all the quantities are equal
$$T_1^\text{eq} = T_2^\text{eq} ...$$
Iterating this process for smaller and smaller partitions of the system proves that,
unsurprisingly, at first order
$$\Delta S = 0.$$

The chemical potential condition
$$\mu^\text{eq}_1 = \mu^\text{eq}_2$$
guarantees chemical equilibrium: supposing an isolated system with 
$\delta Q = \delta W = 0$, $\Delta S>0$ and
$$\Delta S \approx - \left(\frac{\mu_1}{T_1} - \frac{\mu_2}{T}\right)\Delta N_1>0$$
so if $\mu_1>\mu_2$ it must be that
$$\Delta N_1 < 0,$$
in other words, matter flows from higher to lower $\mu$.

The quantity
$$-\gradient\left(\frac{\mu}{T}\right)$$
is a generalized force (see first lesson), much like
$\gradient \left(\frac{1}{T}\right)$.
